\chapter{故障排查与救援}
\chapterlead{排查是按层排除，不是猜。本章按“能不能开机、能不能登录、服务是否运行、应用是否正常”四级给出固定顺序，并覆盖三个最常见事故：磁盘满、内存耗尽、启动失败。}

\section{起步：四级定位}

\begin{booktable}{\textwidth}{J{0.741} J{1.259}}
\toprule
层级 & 判断依据 \\
\midrule
固件与引导 & 有画面、能进引导菜单 \\
内核与 initrd & 内核日志出现，根文件系统挂载成功 \\
init 与网络 & 能登录，SSH 可达 \\
服务与应用 & 单元运行、端口监听、请求有响应 \\
\bottomrule
\end{booktable}

每一级失败，处理方法不同。往下走之前先确认上一级正常，否则会在错误的方向上花时间。

\section{磁盘满}

最常见的线上事故，表现是服务突然写不进去、数据库拒绝启动。

\begin{shellsession}
$ df -h
$ sudo du -xh / --max-depth=1 | sort -h | tail -8
$ sudo find / -xdev -type f -size +500M -exec ls -lh {} + 2>/dev/null | head
$ journalctl --disk-usage
$ nix-collect-garbage --dry-run
\end{shellsession}

四项分别定位：哪个分区满、哪个目录占得多、有没有大文件、日志占多少。NixOS 上还有专门的来源：旧代与 store。

\begin{shellsession}
$ sudo nix-collect-garbage -d
$ sudo nix-store --optimise
\end{shellsession}

清理后仍满，检查被删除但仍被进程占用的文件：

\begin{shellsession}
$ sudo lsof +L1 | head
\end{shellsession}

\texttt{+L1} 列出链接数为零的已打开文件。它们占着空间却看不见，重启对应服务即可释放。

\section{内存耗尽}

表现是系统卡顿、OOM killer 杀掉进程、SSH 连不上。

\begin{shellsession}
$ free -h
$ systemd-cgtop -m --iterations=1 | head -10
$ journalctl -k | grep -i "out of memory" | tail -5
$ cat /proc/pressure/memory
\end{shellsession}

三条命令分工不同：\texttt{systemd-cgtop -m} 按服务显示内存占用，比进程列表更接近实际情况；\texttt{journalctl -k} 中的 OOM 记录会写明被杀的进程与当时的可用内存；\texttt{/proc/pressure/memory} 给出内存压力指标，持续偏高说明已经长期吃紧。

处理分两步：短期限制或重启占用最多的服务；长期加交换空间、限制服务内存（\texttt{MemoryMax}）、或增加内存。

\begin{nixcode}
systemd.services.api.serviceConfig = {
  MemoryMax = "1G";
  MemoryHigh = "768M";
};
\end{nixcode}

\section{启动失败}

系统起不来时，第一步是判断卡在哪一级。

\begin{shellsession}
# 从安装介质启动后
$ sudo mount /dev/disk/by-label/nixos /mnt
$ sudo mount /dev/disk/by-label/boot /mnt/boot
$ sudo nixos-enter --root /mnt
$ journalctl -b -1 -p err
$ nixos-rebuild switch --rollback
\end{shellsession}

\texttt{nixos-enter} 进入目标系统环境，可以在其中查看上一次启动的日志、修改配置、执行回滚。这条路径不需要重新安装。

引导菜单里的旧代是更快的选择：开机时选 Older version 进入上一代，登录后修复配置。

\begin{warnbox}[先确认不是硬件]
反复在同一个位置失败、且旧代也起不来，可能是磁盘或内存故障。用 \texttt{smartctl -a} 查看磁盘健康，用内存测试工具检查内存。软件层面的回滚对硬件故障无效。
\end{warnbox}

\section{网络与 DNS}

\begin{shellsession}
$ ip -brief address
$ ip route
$ ping -c 2 1.1.1.1
$ dig +short example.com
$ resolvectl status | head -8
\end{shellsession}

五条命令依次排除：有没有地址、有没有默认路由、能不能到公网、DNS 是否解析、系统用的是哪个 DNS。能用 IP 访问但域名不行，问题在 DNS；连 IP 都不通，问题在路由或链路。

\section{服务起不来}

\begin{shellsession}
$ systemctl status api
$ journalctl -u api -n 50 --no-pager
$ systemctl cat api
$ sudo systemctl restart api
\end{shellsession}

\texttt{status} 给出退出码与最近日志，\texttt{cat} 显示实际生效的单元内容。常见原因有四种：路径写错、权限不足、依赖的目录或端口被占用、环境变量缺失。四种都能在日志里找到直接证据。

\section{写好一份事故记录}

排查结束后立刻记录，趁记忆还新。

\begin{textcode}
时间：2026-11-25 14:02 至 14:31
现象：api.example.com 返回 502，持续 29 分钟
定位：df -h 显示 / 使用率 100%；du 定位到 /var/log/journal
原因：日志保留期未设上限，单日写入 8GB
处置：journalctl --vacuum-size=1G；配置 SystemMaxUse=1G
后续：为磁盘使用率加告警，阈值 85%
\end{textcode}

六项内容缺一不可：时间、现象、定位过程、根因、处置、后续动作。第五项以外的四项决定了同类问题下次能否更快解决。

\section{练习}

\begin{enumerate}
  \item 用 \texttt{df} 与 \texttt{du} 找出本机占用最大的目录，记录前三个。
  \item 人为写满一个小分区（可用 loop 设备），观察服务报错，再清理恢复。
  \item 用 \texttt{lsof +L1} 找出已删除但仍被占用的文件，说明如何释放空间。
  \item 用 \texttt{nixos-enter} 从上一次启动的日志中找出三条错误。
  \item 按六项格式写一份事故记录，描述你刚制造的磁盘满事件。
\end{enumerate}

\begin{checkpoint}
本章结束时，应能：
\begin{itemize}
  \item 按四级顺序判断故障在哪一层；
  \item 用四项命令定位磁盘满，并知道 NixOS 特有的两个来源；
  \item 用 \texttt{systemd-cgtop} 与 OOM 日志处理内存问题；
  \item 用引导菜单与 \texttt{nixos-enter} 修复启动失败；
  \item 按六项格式写出一份可复用的记录。
\end{itemize}
\end{checkpoint}
